% ===================================================================
%  TABELLA N-RO 7 — Le village durante le hiberno. — Le casa rural
%  durante le primavera.
%  Traduction 2026 d'apres texto/io, controlee sur texto/fr et
%  texto/es. Consigne : voir texto/ia/10-tabella-01.tex.
% ===================================================================

%%K t07-noto-1 noto
\VUnotes{143.2mm}{%
(*) Vide quanto al detalios del repastos, le tabellas 6 e 13.%
}

%%K t07-apar-1 apar
\VUsaut{4.0mm}
\VUcentre{13.2pt}{90}{SECUNDE SERIE}
\VUsaut{3.0mm}
\VUfilet{16mm}
\VUsaut{5.6mm}
\VUtitre{15.0pt}{60}{TABELLA N\textsuperscript{o} 7}
\VUsaut{3.0mm}

%%K t07-tit-1 sub
\VUcentre{12.0pt}{0}{\textit{Prime scena.}}
\VUsaut{3.0mm}

%%K t07-tit-2 sub
\VUpk{10.2pt}{40}{Le village durante le hiberno.}
\VUsaut{4.0mm}

%%K t07-c1-01-1 p
1. --- Durante le vacantias del nove anno, io accompaniava mi patre a Villanova,
un parve village ubi ille debeva tractar alcun affaires commercial.

%%K t07-c1-02-1 p
2. --- Nos usava le \VUgras{diligentia} \textsuperscript{(11)} que functiona
inter le citate e ille burgo; ma proque il faceva multo frigido, iste viage in
un vehiculo pauco confortabile non esseva multo agradabile.

%%K t07-c1-03-1 p
3. --- Pro isto nos sentiva un ver placer quando nos videva le \VUgras{cochero}
haltar su vehiculo sur le placia, ante le \VUgras{albergo}
\textsuperscript{(2)} del \textit{Urso Blanc}. \VUgras{Le albergero}
\textsuperscript{(4)} nos attendeva sur le limine de su casa, fumante
tranquillemente su \VUgras{pipa}.

%%K t07-c1-03-2 p
Ille nos invitava entrar e io non me faceva rogar pro installar me ante le foco
de sarmentos que crepitava in le focar. Alora io me burlava multo del acute
\VUgras{vento del nord} que face grunnir le vetule \VUgras{insignia}
\textsuperscript{(3)} e portava via in nigre turbines \VUgras{le fumo}
\textsuperscript{(1)} exiente del camino.

%%K t07-c1-04-1 p
4. --- Il esseva le hora del dejuner. Sin attender plus longe, nos mangiava con
appetito le simple ma sapide repasto que on nos serviva (*).

%%K t07-tit-3 sub
\VUpk{10.2pt}{40}{Alcun visitas.}
\VUsaut{3.0mm}

%%K t07-c1-05-1 p
5. --- Post le dejuner, io accompaniava mi patre, qui es agente de assecurantias,
in su clientela. Primo nos visitava le \VUgras{ferrero} \textsuperscript{(5)},
qui habita presso le albergo. Ille esseva occupate a ferrar un cavallo. Io
admirava le robustessa de su companion qui teneva solidemente sur su avantal de
corio le \VUgras{ungula} del cavallo, durante que le ferrero fixava le
\VUgras{ferro} \textsuperscript{(6)} per forte colpos de martello.

%%K t07-c1-06-1 p
6. --- Postea nos iva al \VUgras{scarpero} \textsuperscript{(7)} del village, le
vetule Jacobo. Ben que ille ha como insignia un belle \VUgras{botta}, ille se
contentava de resolar un vetule par de scarpas traforate.

%%K t07-c1-06-2 p
Le ancian nos diceva ridente : « In le citate, io esseva “fabricante de
scarpas” e io non habeva clientes. Hic, io es “reparator de scarpas” e le labor
es multe. »

%%K t07-c1-07-1 p
7. --- Ille nos parlava de su vicino le \VUgras{barbero} \textsuperscript{(8)},
del qual le clientela non es contente. Su \VUgras{rasorios} es sempre brecate e
su \VUgras{cisorios} nunquam talia ben. Ma proque ille es le sol barbero del
village, on es obligate naturalmente a facer se rader e facer taliar su capillos
per ille. Del resto ille es un homine avar e disagradabile.

%%K t07-c1-08-1 p
8. --- Felicemente, le altere \VUgras{vicinos} non le resimila. Per exemplo le
\VUgras{carrettero} \textsuperscript{(9)} es un homine benigne, laboriose e
complacente. Il es un placer facer reparar un vehiculo per ille. Su labor es
sempre finite.

%%K t07-c1-09-1 p
9. --- Le vetule Jacobo anque non se plangeva del \VUgras{sellero}
\textsuperscript{(10)}, le qual ille alcun vices adjuta a suer le
\VUgras{harnesamentos} quando le labor urge.

%%K t07-c1-09-2 p
Mi patre le questionava super su familia : « Totes sta ben. Mi nepta es in le
\VUgras{schola} \textsuperscript{(12)}. Su \VUgras{maestra}
\textsuperscript{(13)} es multo contente de illa, illa es un bon \VUgras{alumna}
\textsuperscript{(14)}. Illa non resimila mi mal puero de filio; ille pensa
solmente a jugar. Le \VUgras{maestro} \textsuperscript{(15)} non succede
inseniar le lo que sia. »

%%K t07-tit-4 sub
\VUsaut{3.0mm}
\VUpk{10.2pt}{40}{Conversationes de village.}
\VUsaut{3.0mm}

%%K t07-c1-10-1 p
10. --- Le ancian postea nos racontava le novas del village. On va construer un
nove schola, proque illo installate in le \VUgras{casa communal} deveniva troppo
parve. Del resto isto es facile, proque il bastara comprar un parte del jardin
del \VUgras{molinero}. Isto essera multo profitabile pro le povre homine. Pro le
frigido, le \VUgras{molas} del \VUgras{molino} \textsuperscript{(16)} non pote
plus moler le tritico, proque le \VUgras{rota} \textsuperscript{(17)} esseva
haltate per le gelo del \VUgras{glacie} \textsuperscript{(25)} del
\VUgras{rivetto} \textsuperscript{(23)}.

%%K t07-c1-11-1 p
11. --- « Quanto al constructiones, ha vos vidite nostre nove \VUgras{ecclesia}
\textsuperscript{(18)}? Illo es multo pittoresc sub su mantello de nive. Le
\VUgras{corvos} \textsuperscript{(19)}, qui non recognosce plus lor vetule
campanario, tristemente se posa sur le grande arbore del \VUgras{cemeterio}
\textsuperscript{(20)} ».

%%K t07-c1-12-1 p
12. --- Iste idea del cemeterio rememorava al scarpero le personas que mi patre
habeva cognoscite e qui moriva le anno passate. « Quante \VUgras{tumbas}
\textsuperscript{(21)}, mi bon senior, se adjungeva al alteres, sub le
\VUgras{cypresso} \textsuperscript{(22)}! Le morte falceva nostre vetule
notario. Tote le villageses assisteva a su \VUgras{funerales}. »

%%K t07-c1-13-1 p
13. --- « Ecce su successor qui patina sur le rivetto. Ille justo veniva facer
reparar per me le corregia de su \VUgras{patin} \textsuperscript{(26)} que
esseva rumpite. »

%%K t07-c1-14-1 p
14. --- « Ecce anque sur le ripa del rivetto, le nove \VUgras{guarda rural}
\textsuperscript{(27)}. Ille es un ex-gendarme qui terrifica le mal pueros del
village. Illes fugi tosto que illes vide su \VUgras{gambieras}
\textsuperscript{(29)} e su \VUgras{kepi} \textsuperscript{(28)}. »

%%K t07-c1-15-1 p
15. --- « Ha! ecce le vetule Joseph qui conduce un \VUgras{carretta de fasces de
ligno} \textsuperscript{(30)} pro le panetero. --- Isto es ver, diceva mi patre,
io recognosce su attelage de \VUgras{mulas} \textsuperscript{(31)}. Io ha ben
besonio de parlar le, io profitara le occasion. A revider, patre Jacobo. »

%%K t07-c1-16-1 p
16. --- Justo quando nos exiva, le pueros del village lassava le schola e
currente se precipitava al \VUgras{glissatorio} \textsuperscript{(33)} con le
risco de cader e rumper se un membro in le \VUgras{cadita}
\textsuperscript{(34)}. Un de illes prendeva su \VUgras{traino}
\textsuperscript{(32)} presso le carrettero, faceva seder in illo un de su
camerados, e partiva currente.

%%K t07-c1-17-1 p
17. --- Alteres se precipitava a un \VUgras{amasso de nive}
\textsuperscript{(37)} presso le \VUgras{ponte} \textsuperscript{(40)} e
comenciava bombardar le \VUgras{statua de nive} \textsuperscript{(39)} que illes
habeva erigite le die de heri e al qual illes habeva ponite un cappello, un pipa
e un sacco de schola in bandoliera.

%%K t07-c1-18-1 p
18. --- Illes se cura pauco del vento gelide e del frigido. Le majoritate mesmo
non ha un \VUgras{grande cravata} \textsuperscript{(35)} e, pro recalefacer se
post le combatto per \VUgras{bollas de nive} \textsuperscript{(38)}, il les
basta un mano plen de multo calide \VUgras{castanias} \textsuperscript{(42)}
comprate del vetule \VUgras{venditrice} Jacobina, qui treme de frigido sub su
\VUgras{mal chal} \textsuperscript{(43)} troppo tenue.

%%K t07-tit-5 sub
\VUsaut{3.0mm}
\VUcentre{12.0pt}{0}{\textit{Secunde scena.}}
\VUsaut{3.0mm}

%%K t07-tit-6 sub
\VUpk{10.2pt}{40}{Le casa rural durante le primavera.}
\VUsaut{4.0mm}

%%K t07-c2-01-1 p
1. --- Malgrado tote le placeres del hiberno, nos con profunde gaudio saluta le
retorno del \VUgras{primavera}.

%%K t07-c2-01-2 p
Quanto un allegrante tabella pare un corte de ferma, le vespere, in le mense de
maio!

%%K t07-c2-02-1 p
2. --- Al horizonte le \VUgras{sol} \textsuperscript{(1)} lentemente se corica,
inundante le \VUgras{campania} \textsuperscript{(3)} de su purpuree
\VUgras{radios} \textsuperscript{(2)}. Le diligente \VUgras{aratero}
\textsuperscript{(6)} hasta arar su \VUgras{campo} \textsuperscript{(4)}.

%%K t07-c2-02-2 p
Con su \VUgras{aratro} \textsuperscript{(5)} attelate a un vigorose
\VUgras{cavallo} \textsuperscript{(7)}, ille tracia le \VUgras{sulco}
\textsuperscript{(8)} in le qual le \VUgras{seminator} \textsuperscript{(9)} con
un mano plen jectara le \VUgras{semine} que producera le auree spicas.

%%K t07-c2-03-1 p
3. --- Le \VUgras{falcatores} \textsuperscript{(11)} providite de lor falces
falcava le herba del prato, le fenatores siccava le \VUgras{feno}
\textsuperscript{(10)} tornante lo con \VUgras{furcas} \textsuperscript{(12)} e
rastros, e ecce que illes reveni.

%%K t07-c2-03-2 p
Unes marcha ante le pesante carro tirate per boves; illes porta lor utensile sur
le spatula. Alteres, plus indolente, se installava confortabilemente sur le
odorante feno.

%%K t07-c2-04-1 p
4. --- Passante illes face un amical saluto al \VUgras{jardinero}
\textsuperscript{(16)} occupate in le \VUgras{verdiero} \textsuperscript{(13)} a
taliar le \VUgras{arbores fructifere} e graffar le \VUgras{pirieros}
\textsuperscript{(14)}. Le \VUgras{prunieros} \textsuperscript{(15)} comencia a
florer, e, in le \VUgras{jardin de legumines} \textsuperscript{(17)} ben
fertilisate, le legumines, caules e salatas ja sta ben.

%%K t07-c2-04-2 p
Sur le mureto, un belle \VUgras{pavon} \textsuperscript{(18)} disploia su cauda,
pro facer admirar su plumas belle.

%%K t07-tit-7 sub
\VUpk{10.2pt}{40}{Corte de ferma.}
\VUsaut{3.0mm}

%%K t07-c2-05-1 p
5. --- Le portas del \VUgras{stabulo} \textsuperscript{(19)} es aperte e on vide
le mangiatoria e le \VUgras{lecto de palea} \textsuperscript{(23)} del
\VUgras{cavalla} \textsuperscript{(21)} que le \VUgras{palafrenero}
\textsuperscript{(20)} justo disattelava. Le \VUgras{pullo de cavallo}
\textsuperscript{(22)} tanto comic con su longe gambas, seque su matre
saltettante.

%%K t07-c2-06-1 p
6. --- Presso le \VUgras{puteo} \textsuperscript{(29)} le \VUgras{vaccas}
\textsuperscript{(25)} va biber in le \VUgras{gaveta} \textsuperscript{(26)} de
ligno que les servi como abiberatorio. Le \VUgras{vitello}
\textsuperscript{(24)} profita le occasion pro suger, e le \VUgras{tauro}
\textsuperscript{(27)} flairante le bon odor del feniera, mugi gaimente e leva
fiermente le capite con \VUgras{cornos} \textsuperscript{(28)} potente.

%%K t07-c2-07-1 p
7. --- Totes labora. Le \VUgras{servitrice} \textsuperscript{(32)} tene in le
mano le \VUgras{corda} \textsuperscript{(31)} del puteo. Illa hauri aqua con un
\VUgras{situla} \textsuperscript{(30)} pro abiberar le bestial. Presso le
\VUgras{granario} \textsuperscript{(33)}, un homine rastra le paleas e ante le
porta del \VUgras{casa} \textsuperscript{(34)} un vetule \VUgras{filatrice}
\textsuperscript{(35)} con su rota de filar e su \VUgras{conucula}, fila le
\VUgras{lino} o le \VUgras{cannabe} como in le tempores ancian.

%%K t07-tit-8 sub
\VUsaut{3.0mm}
\VUpk{10.2pt}{40}{Le fermero.}
\VUsaut{3.0mm}

%%K t07-c2-08-1 p
8. --- Le \VUgras{fermero} \textsuperscript{(36)} dirige tote iste labores e da
instructiones a su servitores. Ille recommenda mitter sub le tecto le
\VUgras{herpica} \textsuperscript{(40)} e le \VUgras{picco} \textsuperscript{(39)}
que on oblidava presso le \VUgras{cabana} \textsuperscript{(38)} del \VUgras{can}
\textsuperscript{(37)}.

%%K t07-c2-09-1 p
9. --- Su critos espaventa un \VUgras{columba} \textsuperscript{(41)}, que a
tote alas vola in le \VUgras{columbario} \textsuperscript{(42)}, e le
\VUgras{gallinas} \textsuperscript{(43)} qui hasta reentrar le \VUgras{gallinario}
\textsuperscript{(44)}.

%%K t07-c2-10-1 p
10. --- Ante le \VUgras{ovil} \textsuperscript{(48)}, ecce le vetule
\VUgras{pastor} \textsuperscript{(45)} qui reconduce su \VUgras{grege} de oves
\textsuperscript{(49)}. Tenente su \VUgras{baston} de pastor
\textsuperscript{(46)}, que nunquam le manca, como su \VUgras{blusa}
\textsuperscript{(47)} discolorate, ille conduce al stabulo \VUgras{arietes}
\textsuperscript{(50)}, \VUgras{oves} \textsuperscript{(53)}, \VUgras{agnos}
\textsuperscript{(54)}, \VUgras{capras} \textsuperscript{(51)} e
\VUgras{capretos} \textsuperscript{(52)}. Su fidel \VUgras{can}
\textsuperscript{(55)} Argos veni le ultime e incita per su latratos le
diligentia del retardatarios.

%%K t07-c2-11-1 p
11. --- Tote iste ruito e movimento non turba le quietude del bon \VUgras{vacca
lactifere} \textsuperscript{(56)} qui face tintinnar su \VUgras{campanetta}
\textsuperscript{(57)}, durante que on la mulge ante le porta del
\VUgras{stabulo} \textsuperscript{(58)}.

%%K t07-c2-12-1 p
12. Illa pare comprender, que illa debe dar su lacte, pro que le
\VUgras{fermera} \textsuperscript{(60)} pote preparar \VUgras{butyro} in le
\VUgras{butyriera} \textsuperscript{(61)} o le excellente \VUgras{caseos}
\textsuperscript{(64)}, que gutta del tabuletta ponite sur le \VUgras{cuva}
\textsuperscript{(63)}.

%%K t07-c2-13-1 p
13. --- Tote le \VUgras{lacteria} \textsuperscript{(59)} es mantenite multo
nette per le \VUgras{servitor de ferma} \textsuperscript{(65)} qui justo lavava
le \VUgras{pottos a lacte} \textsuperscript{(62)} in le grande situla que ille
porta con le mano.

%%K t07-tit-9 sub
\VUsaut{3.0mm}
\VUpk{10.2pt}{40}{Le corte de aves.}
\VUsaut{3.0mm}

%%K t07-c2-14-1 p
14. --- Con su spisse sabots de ligno, ille marcha un pauco pesantemente, e assi
face fugir le \VUgras{gallo de India} \textsuperscript{(68)}, le \VUgras{anates
femina} \textsuperscript{(66)}, le \VUgras{anates mascule} \textsuperscript{(67)}
e le \VUgras{ansores} \textsuperscript{(69)} qui large aperi \VUgras{lor becco}
\textsuperscript{(70)} e crita inquiete.

%%K t07-c2-14-2 p
Le \VUgras{gallo} \textsuperscript{(71)}, plus battaliose, erige su rubie
\VUgras{cresta} \textsuperscript{(72)}, succute le plumas de su \VUgras{cauda}
\textsuperscript{(73)} e face audir un « cocorico » sonorose.

%%K t07-c2-15-1 p
15. --- Un \VUgras{gallina matre} \textsuperscript{(74)} rapina sur le
\VUgras{fumo} \textsuperscript{(81)} con su \VUgras{parve pullos}
\textsuperscript{(75)}. Le \VUgras{porchetto} \textsuperscript{(78)} fugi verso
su matre, le enorme \VUgras{porca} \textsuperscript{(77)} que nos vide presso le
porchiera.

%%K t07-c2-16-1 p
16. --- In lor compartimentos, le \VUgras{conilios} \textsuperscript{(79)}
mangia avidemente le delicate \VUgras{herba} \textsuperscript{(80)} que les
apporta le filia del fermero. Johannetta ama multo su conilios e, cata die, post
haber cercate le fresc \VUgras{ovos} \textsuperscript{(76)} in le gallinario,
illa veni visitar su parve amicos.
\VUsaut{6.0mm}
\VUfilet{22mm}
